% Volume I: The Epistemology of Suspicion
% Part II. Events, Records, and Reconstructions
% Chapter 6: The Event Is Not the Record
% Spine: proof-spines/ch006-spine.md

\chapter{The Event Is Not the Record}
\label{ch:ch006}


The event is not the record because the ontology of occurrence and preservation already contains several transformations:
\[
X \xrightarrow{\tau} T
\xrightarrow{\kappa} R
\xrightarrow{\mu} M.
\]
Here \(X\) is the event, \(T\) its trace, \(R\) a deliberately or automatically retained record, and \(M\) a representation that has been interpreted, edited, indexed, or narrated. Each arrow can lose information, add structure, or alter what later counts as salient.

The distinctions matter. A trace is a physical consequence of an event; a record is a trace that survives in some preserved form; a representation is a record already shaped for reading, retrieval, or judgment. Preservation therefore does not entail accessibility:
\[
\operatorname{Recorded}(x)
\centernot\Rightarrow
\operatorname{Retrievable}(x)
\centernot\Rightarrow
\operatorname{Recognized}(x).
\]
A fact can be recorded yet buried, retrievable yet unnoticed, and noticed yet misrecognized. This chapter's preface to the rest of Part II is that documentary survival must never be mistaken for transparent presence.
